\section{High-Frequency GaN Switching}
\label{sec:bg-gan}

At \SI{1}{\mega\hertz}, switching loss is sensitive to the energy dissipated per transition:
\begin{equation}
P_{\mathrm{sw}}\approx f_s(E_{\mathrm{on}}+E_{\mathrm{off}}).
\end{equation}
Zero-voltage switching reduces turn-on loss by using circuit current to charge and discharge the half-bridge capacitances before the incoming device conducts. With nonlinear capacitance, the stored energy is
\begin{equation}
E_{\mathrm{oss}}(V)=\int_0^V vC_{\mathrm{oss}}(v)\,\mathrm{d}v.
\end{equation}
The current must also have the correct direction and move the node to the required voltage before the available dead time ends.

Fast GaN transitions increase sensitivity to interconnect inductance, producing \(L\,\mathrm{d}i/\mathrm{d}t\) overshoot, and to capacitive coupling from the switch node into sensing and control. Gate-drive layout, decoupling, commutation-loop area, and return paths therefore affect switching behavior.

Dead time must complete node commutation without excessive reverse conduction. An insufficient interval can leave residual turn-on voltage or allow cross conduction; an excessive interval increases loss. The required interval depends on bus voltage, current, capacitance, and tank state, so burst restart is assessed separately from switching within an active burst.

The prototype uses an integrated LMG2652 half bridge. The supplied discrete-GaN and external-driver models differ in capacitance, delay, sensing response, and thermal paths. These models describe the operating principle, but the assembled stage needs its own switching measurements.

ZVS must be checked from switch-node voltage and the incoming gate command or device turn-on time on the same time base, with adequate bandwidth and controlled probe loading, loop area, and common-mode error. Chapter~\ref{ch:results} applies these criteria to the available captures.
